\documentclass[conference]{IEEEtran}
\IEEEoverridecommandlockouts

\usepackage{cite}
\usepackage{amsmath,amssymb,amsfonts}
\usepackage{graphicx}
\usepackage{textcomp}
\usepackage{xcolor}
\usepackage{url}

\begin{document}

\title{Multi Omics Integration}

\author{
\IEEEauthorblockN{Aman Kumar Singh}
\IEEEauthorblockA{
SEMTM0045 – Data Science Methods and Practice\\
University of Bristol}
}

\maketitle

\begin{abstract}
Provide a concise overview of the research problem, the identified gap, the proposed methodological direction, and the expected impact of the project.

\end{abstract}

\begin{IEEEkeywords}
keyword1, keyword2, ..
\end{IEEEkeywords}

\section{Introduction}
\footnote{GitHub link: \url{https://github.com/UoB-DSMP2026/research-portfolio-fp25098}}



\section{Literature Review}

Multi-omics integration has importantly advanced personalised medicine by providing a broader understanding of complex biological systems beyond single omics analyzes \cite{acharya_comprehensive_2024}. These approaches offer deeper insight into disease mechanisms and patient heterogeneity by integrating diverse modalities such as genomics, transcriptomics and epigenomics, etc. \cite{hong_multi-omics_2025}. However, despite rapid advancements in the field, it continues to face fundamental computational, methodological, and evaluation challenges—most notably, the curse of dimensionality. \cite{guarino_intrinsic-dimension_2023}. Multi-omics datasets are inherently characterised by extremely high-dimensional feature spaces, which are further coupled with relatively small sample sizes and low signal-to-noise ratios\cite{kioroglou_multi-omic_2025,hong_multi-omics_2025}. This high-dimensional, low sample size (HDLSS) setting requires rigorous preprocessing pipelines, which include standardised normalisation, robust batch effect correction, data imputation and effective dimensionality reduction to enable meaningful downstream analysis \cite{abdelaziz_multi-omics_2024}.\\
Even though dimensionality reduction techniques are commonly used to manage high-dimensional multi-omics data, the different methods vary considerably in how well they preserve biological interpretability. In particular, several approaches belong to a category that does not maintain explicit feature-level representations.\cite{shakyawar_icluf_nodate, hong_multi-omics_2025, yang_mdicc_2022}. For instance, Yang et al. perform dimensionality reduction implicitly through MDICC \cite{yang_mdicc_2022} via low-rank subspace learning without constructing an explicit latent feature space, which limits biological interpretability despite strong clustering performance. Similarly, Shakyawar et al. \cite{shakyawar_icluf_nodate} reduce dimensional complexity through iCluF by iterative refinement of patient–patient similarity network, yet provide limited insight into the underlying omics features contributing to the learned representations. Taken together, these methods highlight a broader trend: gains in dimensionality reduction performance and computational efficiency often come at the expense of transparency and biological interpretability.\\
In contrast, classical approaches such as PCA  and NMF provide explicit feature contributions which enable clearer biological interpretation \cite{downing_primer_nodate}. Network-based methods such as Similarity Network Fusion (SNF)\cite{cantini_benchmarking_2021} integrate multi-omics through fused similarity networks and factor-based models such as MOFA \cite{darkaoui_integrative_2026} learn latent factors that retain partial biological interpretability. However, these methods are often limited in capturing complex nonlinear relationships inherent in biological systems. This comparison highlights a fundamental trade-off in the literature between interpretability and modelling expressiveness with recent methods increasingly favouring predictive performance over biological transparency.\\
Batch effect correction represents an equally critical yet frequently underemphasised preprocessing step beyond dimensionality reduction. Batch effects, which arise from variations in experimental conditions, sequencing platforms or laboratory protocols, can obscure true biological signals and introduce systematic bias into downstream analysis\cite{goh_why_2017}. Consequently, even well-designed dimensionality reduction techniques may fail when applied to uncorrected data. In spite of this, it has been observed that several studies disregarded robust batch correction strategies\cite{shakyawar_icluf_nodate, yang_mdicc_2022, hira_integrated_2021}. For example, Hira et al. \cite{hira_integrated_2021} combine multi-source omics data (TCGA and GDC) without explicit batch correction, and even though normalisation and sample intersection are applied, their direct concatenation likely preserves technical irregularity, which leads to representations influenced by batch-driven artefacts rather than biological signals.
Similarly, Yang et al. \cite{yang_mdicc_2022} focus on data harmonisation using MDICC but do not explicitly address cross-platform batch effects.
On the other hand, Zhang et al. \cite{zhang_integrative_2025} use ComBat, whereas Kioroglou et al. \cite{kioroglou_multi-omic_2025} use covariate filtering, which demonstrates improved control of non-biological variation. The omission of such steps constitutes a methodological gap because uncorrected technical noise propagates through the analytical pipeline and compromises downstream validity.\\
Within the downstream analysis phase, deep learning models represent a prominent methodological archetype frequently deployed for multi-omics integration tasks. Architectures such as ANNs, GNNs, GCNs and transformer-based models are widely deployed to capture complex nonlinear relationships and often outperform traditional machine learning methods in predictive tasks\cite{sartori_comprehensive_2025}. But in spite of these performance gains, these models remain inherently black-box, which offers limited feature-level interpretability and biological insight, which further restricts their clinical applicability \cite{kang_roadmap_2022}. For example, Wu et al. \cite{wu_deepmoic_2024} demonstrate better performance with DeepMoIC over classical classifiers such as SVM, Random Forest and KNN, but fail to provide clear attribution of the molecular features driving the model's predictions, which further limits its biological relevance. Post-hoc Explainable AI (XAI) frameworks like SHAP, LIME, and Integrated Gradients have been introduced to extract feature-level importance from such complex models. However, their use in multi-omics studies remains limited. While these Explainable AI frameworks can support biomarker prediction and enhance clinical applicability, they primarily provide approximate explanations and do not fully overcome the inherent opacity of deep learning models. \cite{noauthor_full_nodate}.\\
This limitation is further magnified by the HDLSS nature of multi-omics data. Deep learning models, which are inherently data-intensive, are prone to overfitting in such settings, which leads to the decreasing generalisability \cite{kang_roadmap_2022}. In comparison, traditional machine learning methods—particularly ensemble approaches such as XGBoost and Random Forest show greater robustness and give inherent feature importance measures that enhance interpretability. Empirical comparisons reinforce this distinction: Wang et al. \cite{wang_mogonet_2021} show that XGBoost outperforms their proposed deep learning model, MOGONET, on certain datasets and achieves comparable performance on others, wheras Hong et al. \cite{hong_multi-omics_2025} show that LightGBMXT achieves superior predictive performance (AUC = 0.9727) compared to multiple deep learning architectures. These findings challenge the implicit assumption that deep learning is universally superior and instead suggest that simpler models may be more appropriate in HDLSS settings.\\
Another persevering and largely unresolved challenge is the lack of generalisability and robustness of proposed models. Even though several studies report high predictive performance still results are often derived from limited validation strategies, which include single-cohort analyses and insufficient cross-validation \cite{shakyawar_icluf_nodate, kioroglou_multi-omic_2025, gillenwater_multi-omics_2021}. For example, Kioroglou et al. demonstrate the potential of multi-omics profiling but restrict their analysis to a single cohort without external validation, which limits its generalisability \cite{kioroglou_multi-omic_2025}. Similarly, Hong et al., in spite of reporting high performance, still rely on a small dataset with a simple hold-out validation strategy, which raises concerns regarding overfitting and optimistic bias \cite{hong_multi-omics_2025}. Moreover, the lack of functional validation of identified biomarkers and limited integration across diverse omics layers undermines reproducibility and biological credibility.\\
More rigorous validation strategies provide partial solutions. Benkirane et al. evaluate their framework across multiple independent datasets, which demonstrate improved robustness \cite{benkirane_customics_2023}. Likewise, Wu et al. validate their framework, DeepMoIC, across multiple TCGA cohorts (e.g., BRCA, LGG, KIPAN) and offer broader evaluation \cite{wu_deepmoic_2024}. However, even such approaches have limitations because validation within a single data repository (e.g., TCGA) may not capture real-world diversification due to shared batch effects, population biases and experimental protocols. Eventually, multi-cohort validation represents progress, but it does not fully resolve the challenge of generalisability.\\
Collectively, these findings reveal a critical disconnect in
current multi-omics research. The advances in integration
and predictive modelling have improved performance, but these advances are not matched by equivalent progress in interpretability, validation and robustness. This imbalance restricts clinical translatability because effective personalised medicine not only requires accurate predictions but also biologically interpretable and reproducible insights. Addressing these interconnected limitations and especially the trade offs between performance, interpretability, and generalisability remains a central challenge
and defines a clear research gap for future work.


\section {Research Gap and Objectives}
Current multi-omic integration approaches exhibit a fundamental imbalance between predictive performance, interpretability, robustness and generalisability in spite of substantial progress in multi-omics integration. As identified in the preceding literature review, many frameworks achieve high accuracy but rely on complex latent representation learning methods that fail to retain explicit feature-level information, thus limiting biological interpretability. This issue is further compounded by inadequate preprocessing strategies, particularly the omission of robust batch effect correction and systematic handling of missing data, which introduces technical bias and compromises downstream validity.\\
Moreover, the high-dimensional, low sample size (HDLSS) nature of multi-omics data poses significant challenges for model robustness. Deep learning-based approaches, while powerful in capturing nonlinear relationships, are prone to overfitting and often lack transparency, restricting their clinical applicability. In parallel, although traditional machine learning methods offer improved interpretability and robustness, they are not consistently leveraged within integrative frameworks.\\
A further limitation lies in the evaluation of existing models, where performance is frequently assessed using limited validation strategies, such as single-cohort analyses or simple hold-out splits. The lack of rigorous external validation across heterogeneous datasets raises concerns regarding generalisability and robustness. Consequently, there exists a clear need for a unified multi-omics integration framework that systematically addresses preprocessing deficiencies, balances predictive performance with interpretability, and ensures robust validation for clinically meaningful biomarker discovery.\\
\textbf{Research Question:} 
How can a pan-cancer multi-omics integration framework be designed to improve predictive performance, interpretability, and robustness using standardised preprocessing and explainable models, and be validated through cross-validation and external cohort testing for biomarker discovery?

\textbf{Objectives:}
\begin{itemize}
    \item Develop a robust preprocessing pipeline incorporating standardisation, missing data imputation (e.g. KNN, multiple imputation), and batch effect correction (e.g., ComBat, Harmony) to reduce technical variability across multi-omics datasets.
    \item Implement interpretable dimensionality reduction and integration methods (e.g. PCA, MOFA, NMF) that preserve feature-level contributions while reducing data complexity.
    \item Design a predictive multi-omics modelling framework integrating both deep learning and traditional machine learning models (e.g., neural networks, random forest, XGBoost), alongside explainability methods (e.g., SHAP), to enhance interpretability and identify biologically meaningful features.
    \item Evaluate the framework using rigorous validation strategies, including cross-validation and external validation on independent heterogeneous cohorts to ensure robustness, generalisability, and clinical applicability.
\end{itemize}

\section{Proposed Methodology}
To address the imbalance between predictive performance, feature-level interpretability, and cross-platform generalisability in multi-omics research, this study proposes an integrated analytical framework. The framework begins with rigorous preprocessing, including standardisation, missing data imputation, and batch effect correction, to reduce technical variability. It then applies interpretable dimensionality reduction and integration methods to reduce complexity while preserving biologically meaningful features. Predictive models are subsequently developed using robust machine learning  and deep learning approaches combined with explainability techniques to identify key molecular drivers. Finally, the framework is evaluated through cross-validation and external validation on independent or heterogeneous cohorts to ensure robustness and generalisability.\\
The \textbf{technical pipeline} of the proposed framework is presented in a stepwise manner below.
\subsubsection{\textbf{Data Collection and Preprocessing}} High-dimensional multi-omics datasets are obtained from globally curated repositories, including TCGA, ICGC, and NCBI GEO, enabling cross-cohort analysis across diverse cancer types. To ensure consistency while preserving cohort-specific characteristics, each dataset is processed independently through parallel preprocessing pipelines. Within each cohort, modality-specific normalisation techniques are applied (e.g., TPM/CPM scaling for RNA-Seq and Noob normalisation for DNA methylation) to harmonise measurement scales. Missing values are subsequently handled using robust non-parametric imputation methods, such as MissForest or KNN imputation, ensuring minimal information loss and maintaining data integrity for downstream analysis.
\subsubsection{\textbf{Leave-One-Cohort-Out (LOCO) Strategy}} To prevent data leakage and ensure a fair evaluation, a Leave-One-Cohort-Out (LOCO) strategy is applied at an early stage of the pipeline. In this approach, one entire dataset (for example, GEO) is set aside as an independent test cohort and is not used during model development. All preprocessing steps—such as batch effect correction, feature selection, and data integration—are performed only on the remaining training cohorts (e.g., TCGA and ICGC). The trained model is then evaluated on the held-out cohort, providing a realistic assessment of how well it generalises to completely unseen data from a different source.
\subsubsection{\textbf{Horizontal Integration and Batch Correction}} To increase sample size and partially mitigate the \( p \gg n \) problem, patient data from multiple cohorts are horizontally integrated by stacking samples across datasets. Following integration, modality-specific batch effect correction is applied to address cross-platform technical variability. \textit{ComBat-seq}, standard \textit{ComBat}, and \textit{limma}-based linear models are used depending on the data modality (sequencing, methylation, and proteomics, respectively). This ensures that batch effects are effectively reduced while preserving biologically meaningful variation across modalities.
\subsubsection{\textbf{Feature Selection and Multi-Omics Integration}} To reduce computational complexity and avoid supervised bias, each modality is independently filtered using unsupervised variance-based criteria, with features ranked via Median Absolute Deviation (MAD) and Coefficient of Variation (CV), retaining only the most variable features for downstream analysis. The resulting filtered datasets are then integrated and compressed using MOFA+ or Joint Non-Negative Matrix Factorisation (Joint NMF), yielding a low-dimensional latent representation that captures shared structure across omics layers.
\subsubsection{\textbf{Multi-Architecture Benchmarking}} The fused latent representation is evaluated using a unified benchmarking framework. L1-regularised logistic regression is used as the baseline model, and compared against classical methods (KNN, linear and RBF SVM), ensemble models (Random Forest, XGBoost, LightGBM), and deep learning models (ANN, GCN, and transformer-based architectures). This enables systematic comparison across increasing model complexity.
\subsubsection{\textbf{Nested Clinical Validation Framework}} The models are evaluated using a nested validation strategy. Internal validation is performed using a Leave-One-Cohort-Out (LOCO) approach to generate robust performance estimates and guide model selection. External validation is subsequently conducted on completely independent cohorts, including datasets from cBioPortal (MSK-IMPACT Clinical Sequencing) and the TARGET project, using inference-only testing to assess true real-world generalisability.
\subsubsection{\textbf{Translational Interpretability and Visual Validation}} To ensure biological interpretability of the learned representations, low-dimensional embeddings are visualised using UMAP and t-SNE to assess clustering and tissue-specific separation. Feature-level explanations are obtained using SHAP-based methods (TreeSHAP for ensemble models and KernelSHAP/DeepSHAP for neural networks). These attributions are mapped to MOFA+ loading matrices to trace key latent factors back to original molecular features for biomarker identification. Finally, pathway analysis using GSEA and KEGG mapping is used to validate biological relevance and alignment with known disease mechanisms.

\textit{\textbf{Justification:}} The proposed framework addresses key limitations in existing multi-omics approaches by improving interpretability, robustness, and generalisability. To reduce the impact of the HDLSS setting, multiple datasets from different repositories are integrated horizontally to increase sample size (n), thereby reducing overfitting and improving generalisability. Unlike latent subspace methods (e.g., MDICC, iCluF) that obscure feature-level structure, MOFA+ and Joint NMF preserve interpretability through explicit feature–factor relationships.
\\
\begin{figure}[h]
    \centering
    \includegraphics[width=1.0\linewidth]{Untitled Diagram.drawio (1).png}
    \caption{Proposed Framework}
    \label{fig:proposed_method}
\end{figure}
In contrast to prior studies with limited preprocessing and validation, the framework incorporates modality-specific batch correction (ComBat-seq, ComBat, limma) and rigorous LOCO-based and external validation to ensure reliability. Additionally, a range of models—from linear to ensemble to deep learning—is systematically evaluated, with SHAP-based explanations and pathway enrichment (GSEA, KEGG) used to ensure both interpretability and biological relevance. Overall, the framework provides a more robust, interpretable, and generalisable alternative to existing baselines.\\
\textit{\textbf{Metrics and Validation Criteria:}} L1-regularised logistic regression is used as the baseline, serving as a sparse and interpretable benchmark to assess the added value of multi-omics integration and more complex models. Validation follows a nested strategy, with internal Leave-One-Cohort-Out (LOCO) to prevent data leakage and external testing on unseen cohorts to ensure robustness and real-world generalisability.\\
Model performance is evaluated using AUC-ROC for class discrimination, F1-score, precision, and recall to account for class imbalance, accuracy for overall performance, and log-loss for probabilistic calibration. Statistical significance is assessed using DeLong’s and Wilcoxon signed-rank tests. The structural quality of learned representations is evaluated using Silhouette score and Davies–Bouldin Index (DBI), while biological validity is confirmed through GSEA and KEGG enrichment with FDR correction, ensuring clinically meaningful and functionally relevant findings.




\section{Data, Ethics, and Practical Considerations}
\subsubsection{Overview of Data Sources and Legal Licensing}
\textbf{TCGA (The Cancer Genome Atlas)}: High-quality, standardised multi-omics cancer dataset with consistent clinical annotation; fully open-access under NIH/GDC Public Domain policy, allowing unrestricted academic reuse and model development; \textbf{ICGC (International Cancer Genome Consortium)}: Large-scale, globally diverse cancer genomics dataset with strong population coverage; open-access under CC BY 4.0-equivalent licensing requiring attribution to the consortium; \textbf{NCBI GEO (Gene Expression Omnibus)}: Heterogeneous repository of transcriptomic and epigenomic datasets with variable preprocessing quality; mixed licensing (Public Domain, CC BY-NC, CC BY-SA) requiring non-commercial and share-alike compliance for certain datasets; \textbf{cBioPortal and TARGET}: Curated, clinically annotated cancer genomics datasets with reliable preprocessing pipelines; open academic access suitable for benchmarking and research use under controlled but non-commercial licensing agreements.
\subsubsection{Ethical Red-Teaming} A key limitation of the proposed framework is in its reliance on TCGA and ICGC datasets, which are predominantly derived from European populations, creating an ancestral generalisation blind spot and increasing the risk of misclassifying population-specific variants as universally relevant cancer biomarkers. This imbalance may reduce the model’s clinical reliability when applied to underrepresented populations. Patient re-identification risk also remains a significant ethical concern because the integration of high-resolution transcriptomic read counts with continuous DNA methylation profiles can produce highly unique molecular signatures that, in principle, could be cross-referenced with external genealogy or insurance databases by malicious actors. Furthermore, despite the application of batch correction methods such as ComBat-seq, residual non-linear batch effects may persist across heterogeneous datasets, leading advanced models (e.g., GCNs and Transformers) to overfit to subtle technical artefacts rather than learning true underlying biological patterns.
\subsubsection{Technical Mitigations}To systematically address the vulnerabilities identified during red-teaming, the proposed framework will implement three key algorithmic safeguards. First, the ancestral generalisation blind spot will be reduced using adversarial de-biasing networks, which penalise reliance on ethnicity-correlated genetic variants and encourage learning of universal oncogenic signals. Second, patient re-identification risk will be mitigated by injecting calibrated differential privacy (DP) noise into feature matrices, creating formal privacy guarantees that prevent reconstruction of individual multi-omics profiles. Finally, residual technical artefacts will be controlled through contrastive SHAP-based confounding audits, which act as a post-hoc screening mechanism to detect and remove features correlated with hospital or batch identifiers rather than true biological signals.
\section{Project plan}
\begin{figure}[h]
    \centering
    \includegraphics[width=1.0\linewidth]{Gantt Chart1.png}
    \caption{Gantt Chart}
    \label{fig:Gantt Chart1}
\end{figure}
The project is organised as a sequential and interdependent pipeline to ensure technical feasibility. It begins with Dataset Collection and Cohort Selection, where multi-omics datasets are gathered from repositories such as TCGA, ICGC, and GEO. This is followed by Data Normalisation and Imputation and Batch Effect Correction, where modality-specific scaling, missing value handling, and correction methods (e.g., ComBat, limma) are applied to remove technical variability and ensure data consistency. The cleaned datasets are then combined in Horizontal Data Integration to increase sample size and mitigate the HDLSS problem, after which Feature Selection is performed using variance-based methods (e.g., MAD, CV) to retain the most informative features.

The refined data is subsequently used in Multi-Omics Integration and Dimensionality Reduction, where methods such as MOFA+ or Joint NMF generate a compact latent representation capturing shared biological structure. This representation feeds into Model Development and Benchmarking, where machine learning and deep learning models are trained and compared. The models are then assessed in Model Validation and Statistical Evaluation using LOCO and external validation strategies alongside performance metrics and statistical tests to ensure robustness. Finally, Interpretability, Biological Validation, and Reporting are conducted using SHAP-based explanations, visualisation techniques, and pathway enrichment analyses (GSEA, KEGG) to derive biologically meaningful insights. Each stage builds on the previous one, with clear dependencies and milestones, ensuring a coherent and achievable workflow within the project timeframe. 
\section{Non-Technical Summary}
Modern cancer research generates vast amounts of complex data from patient tissue samples, including genetic information and cellular changes. Combining these different layers of data—known as multi-omics—is essential for advancing personalised medicine, as it helps clinicians better understand patient differences and tailor treatments accordingly.

However, current computational models face a fundamental trade-off. Many achieve high predictive accuracy but operate as “black boxes,” offering little insight into how their decisions are made. This lack of transparency becomes a serious limitation in clinical settings, where understanding the biological basis of predictions is just as important as the predictions themselves.

A deeper issue lies in how these models are developed and evaluated. Data collected from different hospitals and laboratories often contains subtle variations due to differences in equipment, protocols, and processing methods. These variations introduce technical noise. Instead of learning true cancer biology, models can inadvertently learn these hidden patterns, creating an illusion of strong performance. Since many studies rely on limited or weak validation strategies, this problem often goes unnoticed. As a result, models that appear effective in research settings may fail when applied to new patient populations, limiting their real-world usefulness and potentially leading to wasted time and resources.

To address these challenges, this report proposes a structured and reliable analytical framework. The process begins with a rigorous data preprocessing stage, where technical inconsistencies across datasets are corrected to ensure that the data reflects genuine biological signals rather than experimental noise. The framework then applies dimensionality reduction techniques to simplify the data while preserving meaningful biological relationships, maintaining a balance between complexity and interpretability.

Building on this foundation, the study evaluates a range of artificial intelligence approaches, including both traditional machine learning and deep learning models, to identify the most robust and reliable predictive system. Finally, the selected model is tested on completely independent datasets to ensure that its performance generalises across diverse clinical settings.
By combining transparency, rigorous validation, and methodological balance, this framework aims to produce more trustworthy and clinically applicable results. Ultimately, it enables researchers and clinicians to identify genuine biological markers with greater confidence, supporting more reliable and effective cancer diagnosis and treatment.

\bibliographystyle{IEEEtran}
\bibliography{references}

\end{document}
